\SourceLeaf{460}
\begin{center}\Large\bfseries BEMERKUNGEN ZUM ACHTEN BAND.\end{center}

Der vorliegende achte Band von \textsc{Gauss}' Werken gibt Ergänzungen zu den drei ersten Bänden und zum vierten Band mit Ausnahme des geodätischen Teils. Er enthält die folgenden Beiträge.

Die Abteilungen \textit{Arithmetik, Analysis} und \textit{Numerisches Rechnen} enthalten eine Reihe einzelner Notizen und Aufsätze, die sich nachträglich in \textsc{Gauss}' Nachlass gefunden haben.

Die Abteilung \textit{Wahrscheinlichkeitsrechnung} enthält hauptsächlich eine Anzahl Belege über die Entdeckung der Methode der kleinsten Quadrate.

Die Abteilungen zur \textit{Geometrie} enthalten Notizen und Briefe, die sich auf die Grundlagen der Geometrie, die Geometria situs, Lehrsätze und Aufgaben aus der elementaren Geometrie, die Verwendung komplexer Größen für die Geometrie und die Theorie der krummen Flächen beziehen.

Die Bearbeitung der Abteilungen \textit{Arithmetik und Algebra} und \textit{Analysis und Funktionentheorie} rührt von Herrn \textsc{Fricke} in Braunschweig, die der Abteilungen \textit{Numerisches Rechnen} und \textit{Wahrscheinlichkeitsrechnung} von den Herren \textsc{Börsch} und \textsc{Krüger} in Potsdam und die der \textit{geometrischen} Abteilungen von Herrn \textsc{Stäckel} in Kiel her. Die Notiz \textit{Allgemeinste Auflösung des Problems der Abwickelung der Flächen} (Seite 447 bis 449) ist von Herrn \textsc{Weingarten} in Charlottenburg bearbeitet worden.

Die allgemeine Redaktion wurde von Herrn \textsc{Brendel} in Göttingen besorgt.

\textsc{Schering}, der dem vierten Band einen Ergänzungsband hinzuzufügen beabsichtigte, hat eine Reihe von Vorarbeiten hinterlassen, hauptsächlich Abschriften aus dem Nachlass und dem Briefwechsel, die bei der Herausgabe benutzt werden konnten.\hfill\mbox{\textsc{F. Klein.}}
